% Options for packages loaded elsewhere
\PassOptionsToPackage{unicode}{hyperref}
\PassOptionsToPackage{hyphens}{url}
\documentclass[
  english,
]{article}
\usepackage{xcolor}
\usepackage[margin=1in]{geometry}
\usepackage{amsmath,amssymb}
\setcounter{secnumdepth}{-\maxdimen} % remove section numbering
\usepackage{iftex}
\ifPDFTeX
  \usepackage[T1]{fontenc}
  \usepackage[utf8]{inputenc}
  \usepackage{textcomp} % provide euro and other symbols
\else % if luatex or xetex
  \usepackage{unicode-math} % this also loads fontspec
  \defaultfontfeatures{Scale=MatchLowercase}
  \defaultfontfeatures[\rmfamily]{Ligatures=TeX,Scale=1}
\fi
\usepackage{lmodern}
\ifPDFTeX\else
  % xetex/luatex font selection
\fi
% Use upquote if available, for straight quotes in verbatim environments
\IfFileExists{upquote.sty}{\usepackage{upquote}}{}
\IfFileExists{microtype.sty}{% use microtype if available
  \usepackage[]{microtype}
  \UseMicrotypeSet[protrusion]{basicmath} % disable protrusion for tt fonts
}{}
\makeatletter
\@ifundefined{KOMAClassName}{% if non-KOMA class
  \IfFileExists{parskip.sty}{%
    \usepackage{parskip}
  }{% else
    \setlength{\parindent}{0pt}
    \setlength{\parskip}{6pt plus 2pt minus 1pt}}
}{% if KOMA class
  \KOMAoptions{parskip=half}}
\makeatother
\usepackage{longtable,booktabs,array}
\usepackage{calc} % for calculating minipage widths
% Correct order of tables after \paragraph or \subparagraph
\usepackage{etoolbox}
\makeatletter
\patchcmd\longtable{\par}{\if@noskipsec\mbox{}\fi\par}{}{}
\makeatother
% Allow footnotes in longtable head/foot
\IfFileExists{footnotehyper.sty}{\usepackage{footnotehyper}}{\usepackage{footnote}}
\makesavenoteenv{longtable}
\ifLuaTeX
\usepackage[bidi=basic,shorthands=off,]{babel}
\else
\usepackage[bidi=default,shorthands=off,]{babel}
\fi
\ifLuaTeX
  \usepackage{selnolig} % disable illegal ligatures
\fi
\setlength{\emergencystretch}{3em} % prevent overfull lines
\providecommand{\tightlist}{%
  \setlength{\itemsep}{0pt}\setlength{\parskip}{0pt}}
\usepackage[]{natbib}
\bibliographystyle{plainnat}
\usepackage{bookmark}
\IfFileExists{xurl.sty}{\usepackage{xurl}}{} % add URL line breaks if available
\urlstyle{same}
\hypersetup{
  pdftitle={Hybrid: Time--Frequency and Temporal Normality Modeling for Industrial Time-Series Anomaly Detection},
  pdflang={en},
  hidelinks,
  pdfcreator={LaTeX via pandoc}}

\title{Hybrid: Time--Frequency and Temporal Normality Modeling for
Industrial Time-Series Anomaly Detection}
\author{}
\date{Working draft, 15 November 2026}

\begin{document}
\maketitle

\section{Introduction}\label{introduction}

Industrial monitoring requires more than recognizing abrupt,
high-amplitude departures from a nominal range. A weak deviation can be
significant when it violates a normal response relationship, whereas a
large change can be acceptable when several variables respond coherently
to an operating transition. An anomaly may also persist after its
initial transition: an effective monitor should discover the event and
retain useful evidence during its abnormal steady state. These
requirements are particularly demanding for multivariate streams in
which responses are delayed, operating conditions change, and
target-domain training records are unlabeled rather than independently
certified as healthy.

Time--frequency representations offer a useful way to describe both
local variation and temporal organization. TimeRadar learns a fractional
time--frequency representation and combines reconstruction with
contextual deviation learning \citep{timeradar2026}. LEFT learns
interacting temporal, spectral, and multiscale views, followed by joint
representation learning and reconstruction \citep{left2026}. These
methods motivate a richer representation than a single instantaneous
vector. However, a representation with good reconstruction or prediction
accuracy does not automatically provide a useful anomaly ranking. A
flexible reconstruction model may reconstruct a sustained abnormal
state, and a predictor may improve normal accuracy without increasing
the separation between abnormal responses and difficult normal
transitions. These are empirical risks, rather than universal properties
of either modeling family.

The central question of this work is therefore whether a common learning
and detection criterion can preserve useful level, temporal, and
cross-variable information while distinguishing persistent departures
from legitimate normal-state variation. We study a complete Hybrid model
rather than treating an additional feature branch as sufficient evidence
of improvement. The candidate considered in this draft learns a joint
density of observed causal windows using invertible transformations that
alternate between the time domain and an orthogonal wavelet domain. The
same normalized negative log-likelihood is minimized on normal-reference
training windows and evaluated when the current observation arrives.
This changes the learning and scoring basis; it is not a new weighting
of the previously studied reconstruction and prediction scores.

Several boundaries matter. First, an exact density is not a semantic
definition of normality: a high-likelihood input can still be
undesirable, and a legitimate new operating state can have low
likelihood. Second, a causal endpoint detector may use the current
observation after it arrives, but it must not use later observations.
This is different from predicting the current observation from strictly
earlier inputs. Third, statistical cross-variable dependence is not
evidence of physical causation or a verified directed lag mechanism.
Finally, robustness to unlabeled training contamination requires a
separate experiment after the detector has demonstrated useful baseline
discrimination. Performance that remains unchanged near a detection
floor does not establish robustness.

We evaluate these questions using fixed public-data settings with
complementary purposes. PSM provides a multivariate development
comparison, SWaT supports event and long-duration analyses, and a
run-separated Tennessee Eastman Process (TEP) task provides a controlled
industrial fault setting with a predefined conditional-response subset.
All three have already influenced this research program and are
explicitly treated as development or mechanism evidence. They cannot
establish a new blind-confirmation claim. Strong simple controls,
including historical means and history/conditional Ridge models where
applicable, remain necessary alongside LEFT and TimeRadar.

The intended methodological contribution is a testable
learning--criterion alignment for time--frequency and temporal normality
modeling, not the invention of affine coupling, wavelets, or likelihood
estimation. The experimental contribution must demonstrate a useful
detection effect attributable to that design at an acceptable
false-alarm cost. The present draft does not assert that either
contribution has been established. The completed comparison and
ablations will determine whether the candidate merits independent
confirmation and subsequent multistage adaptation studies.

\section{Related Work}\label{related-work}

\subsection{Time--frequency and multiview anomaly
representations}\label{timefrequency-and-multiview-anomaly-representations}

TimeRadar uses a learnable fractional time--frequency representation
within its foundation-model anomaly-detection framework
\citep{timeradar2026}. Its fractional transform spans different domain
rotations; it should not be described as a frequency-only encoder. LEFT
combines temporal, spectral, and multiscale tokens, selective cross-view
interaction, and an analysis--synthesis reconstruction chain
\citep{left2026}. CATCH instead examines frequency patching with
channel-aware modeling \citep{catch2025}. These works show different
ways of exploiting temporal and spectral structure. Hybrid does not
claim that simply adding a transform or combining view-specific scores
is new. Its current question concerns a joint normality criterion whose
likelihood includes the transformation of the full observed window. The
time-only ablation is required to determine whether the wavelet-domain
transformations provide information beyond the temporal architecture and
parameter count.

\subsection{Temporal dependencies and relationship
anomalies}\label{temporal-dependencies-and-relationship-anomalies}

Forecasting and conditional-response models can expose departures that
are not apparent from individual variables. Long-context forecasting
also motivates temporal encoders, but a forecasting foundation model is
not automatically an anomaly detector; TimesFM, for example, is
presented as a forecasting model \citep{timesfm2024}. In our industrial
relation task, conditional Ridge is a deliberately strong control: any
proposed nonlinear detector must demonstrate value beyond a legal linear
model before a cross-variable representation is credited. Self-lag
denotes dependence on earlier values of the same variable, whereas
cross-lag denotes dependence on earlier values of another variable.
Neither a nonzero routing weight nor an active gradient proves that a
detector identifies changes in those relationships. Mechanism claims
require known relation-preserving and relation-breaking cases, including
marginally similar controls and normal synchronous changes.

CAROTS uses causality-aware augmentations and contrastive learning to
improve multivariate anomaly detection \citep{carots2025}. Its published
method includes a complete augmentation, representation-learning, and
detection procedure, rather than an arbitrary pairwise binary
classifier. Accordingly, relationship-preserving augmentation and
contrastive discrimination cannot be presented as novel merely by
attaching them to a new backbone. The current Hybrid candidate does not
resume the previously stopped relation-routing or isolated
contrastive-loss variants. Instead, it tests whether joint temporal
density modeling yields useful anomaly evidence; any additional claim
about lag-specific detection remains conditional on controlled results.

\subsection{Normality learning, invertible models, and anomaly
criteria}\label{normality-learning-invertible-models-and-anomaly-criteria}

Reconstruction, prediction, prototype distance, and likelihood make
different assumptions about normality. Their training losses can
emphasize different features from those that control event detection.
Real NVP introduced tractable density estimation through invertible
affine transformations with exact log-determinants \citep{realnvp2017}.
Wavelet Flow already combines wavelets with normalizing flows and models
signal scales explicitly \citep{waveletflow2020}. Hybrid reuses these
established principles; a Haar transform inside a flow is not its
novelty. GANF also uses joint density estimation for time-series anomaly
detection through a graph-factorized flow \citep{ganf2022}, and MTGFlow
studies multivariate detection with dynamic graphs and entity-aware
normalizing flows \citep{mtgflow2023}. Thus neither window likelihood
nor graph/dependency modeling is unprecedented.

The current candidate instead tests a particular alternating-coordinate
joint normality model and its sustained-anomaly behavior under fixed
information permissions. An invertible representation does not discard
dimensions through an encoder bottleneck, but this property alone does
not guarantee that the learned density orders anomalies correctly. The
identity-transform Student-t control and the time-only and local-context
ablations test whether learned dependencies and multiscale organization
are useful, rather than attributing improvement to a generic
heavy-tailed score or a change in score units. A publication-level
distinction from existing flow methods would require such empirical
evidence and an appropriate direct comparison; architectural resemblance
alone is not treated as a contribution.

\subsection{Mixed-unlabeled adaptation and evaluation
permissions}\label{mixed-unlabeled-adaptation-and-evaluation-permissions}

Mixed-unlabeled adaptation raises a different question from
normal-reference learning. A target adaptation pool can contain unusual
but legitimate states, anomalous episodes, and distributional changes
relative to the source data. TSAD-C studies contaminated multivariate
anomaly detection using a complete graph-conditional diffusion framework
\citep{tsadc2025}. Such work motivates contamination-aware evaluation,
but does not justify adding a robust loss when an experiment has not
identified contamination-specific harm. We separate original-background
adaptation, anomalous replacement, and matched normal replacement, and
distinguish adaptation effects from calibration effects. A full
contamination-robustness claim requires both detectable anomalies and a
reproducible benefit under matched training exposure.

Information permissions also limit comparisons. A pretrained source
detector, a target-trained detector, and an offline window detector
answer different questions. Restricting an offline score to a common set
of timestamps does not remove the future observations used to construct
that score. We therefore report native offline and causal endpoint
results separately, disclose target-training and standardization
permissions, and use a shared evaluator only where the raw axes and
labels are compatible. Published scores under different permissions are
background references rather than interchangeable rows of a matched
experiment.

\section{Problem Formulation}\label{problem-formulation}

\subsection{Observations, causal windows, and anomaly
events}\label{observations-causal-windows-and-anomaly-events}

Let \(x_t\in\mathbb R^C\) denote the multivariate observation received
at time \(t\), and let \(\mathcal F_t=\sigma(x_1,\ldots,x_t)\) denote
the available observation history. An endpoint detector produces \(s_t\)
measurable with respect to \(\mathcal F_t\). For a context length \(H\),
its input is

\[X_t=[x_{t-H+1},\ldots,x_t]\in\mathbb R^{H\times C}.\]

The endpoint permission permits the current observation \(x_t\) after
arrival and excludes \(x_{t+1:}\). A predictor of \(x_t\) has the
stricter input permission \(\mathcal F_{t-1}\); the two tasks must not
be conflated. The present density model is an observed-window endpoint
detector, not a masked-target predictor. Window normalization and every
model statistic must respect the same permission.

Ground-truth anomaly events are maximal contiguous positive intervals in
the fixed official label sequence, subject to the preregistered boundary
and censoring rules. Their labels are used only for evaluation after
model and calibration choices have been frozen. A weak anomaly refers to
a small change relative to an explicitly defined training scale, not an
uncertain label or a short event. A sustained anomaly is assessed
through its event-level discovery and within-event coverage, including a
predefined steady portion when one exists. Public fault labels alone do
not identify which variable, lag, or physical relationship is abnormal.

\subsection{A common normality-learning and scoring
criterion}\label{a-common-normality-learning-and-scoring-criterion}

Channelwise location and scale \((\mu_c,\sigma_c)\) are estimated using
legal FIT data only. With a fixed positive numerical floor on scale, the
model receives \(\widetilde X_{t,c}=(X_{t,c}-\mu_c)/\sigma_c\). This
operation preserves within-window level relative to FIT; it does not
subtract each window's own mean. Let
\(f_\theta:\mathbb R^{H\times C}\rightarrow\mathbb R^{H\times C}\) be
invertible, and write \(z_t=f_\theta(\widetilde X_t)\). With a
factorized Student-t prior of fixed degrees of freedom \(\nu=5\),

\[p_0(z)=\prod_{i=1}^{HC} \operatorname{StudentT}_{5}(z_i),\]

the window negative log-likelihood in standardized coordinates is

\[\ell_\theta(\widetilde X_t)=-\log p_0(f_\theta(\widetilde X_t))-\log\left|\det J_{f_\theta}(\widetilde X_t)\right|.\]

The fixed per-domain standardization contributes a constant density
offset in physical coordinates. Its omission from standardized-space
scores does not change within-domain ranking or an identically
calibrated threshold, but raw likelihood values from separately
standardized domains are not assumed comparable. The training objective
and endpoint score are

\[\mathcal L_{\mathrm{FIT}}(\theta)=\frac{1}{|\mathcal T_{\mathrm{FIT}}|}\sum_{t\in\mathcal T_{\mathrm{FIT}}}\frac{\ell_\theta(\widetilde X_t)}{HC},\qquad s_t=\frac{\ell_\theta(\widetilde X_t)}{HC}.\]

The model is selected by the same normal SELECT likelihood, without
QUERY labels. No reconstruction, prediction, reference-retrieval, or
temporal accumulation score is added to \(s_t\).

The current implementation design uses four affine coupling blocks,
alternating time-domain coordinates and a two-level orthogonal Haar
transform. A coupling keeps one partition unchanged and applies an
invertible affine map to its complement:

\[u_a=v_a,\qquad u_b=v_b\odot\exp a_\theta(v_a)+b_\theta(v_a),\]

whose log-determinant is the sum of the log-scales \(a_\theta\).
Alternating checkerboard partitions expose different time--channel
coordinates across blocks. The scale and translation functions use
dilated temporal convolutions and cross-channel mixing. An orthogonal
transform has unit absolute determinant, so it does not create an
unaccounted likelihood term. The implemented two-level Haar transform
requires \(H\) divisible by four and uses exact adjacent-pair analysis
and synthesis without padding. This finite observed-window map can mix
coordinates within the available window; it is not a claim that every
latent coordinate has a directed physical causal interpretation.
Coordinatewise terms used to sum the prior and Jacobian are diagnostic
allocations; after cross-channel transformations they are not
per-channel marginal or conditional densities. Only the complete sum
over all \(H C\) coordinates defines the stated joint likelihood.

This construction differs from a fixed normal center because normality
is learned as a transformed distribution rather than distance to a
single reference. It differs from nearest-neighbor retrieval because no
QUERY-dependent memory search defines the criterion. It differs from the
stopped shared prediction--reconstruction design because the training
and detection quantities are the same likelihood. Nevertheless,
whole-window averaging may dilute a weak local departure or delay a
warning, and maximum likelihood may allocate probability poorly under a
new legitimate operating state. These are explicit falsifiable
limitations, not effects removed by the architecture definition.

\subsection{Training regimes and future adaptation
questions}\label{training-regimes-and-future-adaptation-questions}

We distinguish source training (Stage 1), source-calibrated zero-shot
detection (Z-S), target-calibrated detection without parameter updates
(Z-T), unlabeled target adaptation (U), and few-event supervision (F).
Target CAL may define statistics and thresholds only within its
authorized role; it never supplies gradients. Target FIT supplies U
updates without anomaly labels. F supervision is counted in complete
independent support events, with support, validation, and query events
separated. A change from Z-S to target-trained learning changes the
scientific question and is reported as such.

This first evaluation establishes only the candidate's basic detection
capability under its stated training permissions. It does not establish
multisource transfer, a contamination-robust U mechanism, or the value
of few-event labels. Those questions remain in the research framework
and require separate matched tests once the basic criterion is
supported.

\section{Experimental Setup}\label{experimental-setup}

\subsection{Fixed public-data
settings}\label{fixed-public-data-settings}

The data roles are inherited from the existing frozen assets; they are
not reselected using candidate scores. Every input window must be
contained in a legal role or its explicitly authorized causal context.
No physical record is split into overlapping training and query windows,
and no discontinuous records are concatenated. Existing guards are
retained where they exceed the candidate's actual footprint. The
candidate uses \(H=192\) for PSM and SWaT and \(H=100\) for TEP. These
settings are fixed before candidate QUERY results are opened.

\begin{longtable}[]{@{}
  >{\raggedright\arraybackslash}p{(\linewidth - 4\tabcolsep) * \real{0.3333}}
  >{\raggedright\arraybackslash}p{(\linewidth - 4\tabcolsep) * \real{0.3333}}
  >{\raggedright\arraybackslash}p{(\linewidth - 4\tabcolsep) * \real{0.3333}}@{}}
\toprule\noalign{}
\begin{minipage}[b]{\linewidth}\raggedright
Setting
\end{minipage} & \begin{minipage}[b]{\linewidth}\raggedright
Fixed training and selection assets
\end{minipage} & \begin{minipage}[b]{\linewidth}\raggedright
Fixed evaluation scope and interpretation
\end{minipage} \\
\midrule\noalign{}
\endhead
\bottomrule\noalign{}
\endlastfoot
PSM & FIT \([0,79488)\), SELECT \([80032,105984)\), CAL
\([106528,132481)\); 25-channel FIT-only standardization & Existing
complete test axis \([512,87841)\) and 71 complete events; external
comparisons additionally use their exact common scored intersection.
Repeatedly exposed development data. \\
SWaT & Source/normal-reference FIT \([0,296020)\), separate historical
U-FIT \([296564,395237)\), SELECT \([395781,445117)\), CAL
\([445661,495000)\); 31 retained official channels & Existing scored
QUERY \([608535,944919)\) with 22 complete attacks, including nine
longer than 512 samples. The reserved donor interval is not merged into
normal-reference FIT. One operating record, not 22 independent runs. \\
TEP Rieth2017 V1 & Normal simulation runs 1--60 FIT, 61--80 SELECT,
81--100 CAL; run boundaries preserved & Fault runs 201--220, all 20
fault types; the predefined 4/11/14 subset is reported alongside all
types. Existing target axis is samples 101--960, with fault onset at
sample 161. A separate normal-test pool is retained. \\
\end{longtable}

TEP's fixed response task uses XMEAS9 as target and the established
XMV10/XMEAS21 conditional information. The density candidate scores
their complete three-channel joint window; this differs from the
target-XMEAS9 residual scored by the historical regression task. Their
data axes and legal inputs can be matched, but their detection objects
must remain explicit. A gain in joint-window detection alone cannot
establish improved conditional-XMEAS9 modeling. XMV10 is a closed-loop
manipulated variable and is not assumed to be a strictly exogenous
cause. The three-minute sample interval permits conversion of index
delay into elapsed time in this task. Shared normal prefixes across
faults from the same simulation run are not counted repeatedly as
independent false-alarm evidence. Fault-wise prevalence on the scored
axis is high, so AP is interpreted together with ROC-AUC, normal false
alarms, delay, and coverage rather than alone.

SWaT attack labels identify system-level attack intervals; they are not
precise per-channel abnormality labels. Long-event analyses retain the
nine predefined events and separately describe onset, sustained state,
and recovery where the fixed boundaries permit. A low score during an
event does not by itself prove that a sustained observable deviation
exists in the retained channels. PSM, SWaT, and these TEP assets have
historical development exposure. No result from this evaluation is
described as a blind confirmation, and no favorable entity or fault is
substituted after evaluation.

\subsection{Comparisons, information permissions, and
ablations}\label{comparisons-information-permissions-and-ablations}

Hybrid is the proposed method. LEFT, TimeRadar, and task-appropriate
statistical models are external comparators; old Hybrid versions are
internal historical controls. Native LEFT preserves its published
learning and scoring chain. Its causal endpoint readout is named
explicitly when used. Native retrospective LEFT and TimeRadar scores
remain in an offline comparison track if their construction uses future
observations relative to an earlier scored point. Reusing those scores
on a shared axis does not change their information permission.

\begin{longtable}[]{@{}
  >{\raggedright\arraybackslash}p{(\linewidth - 4\tabcolsep) * \real{0.3333}}
  >{\raggedright\arraybackslash}p{(\linewidth - 4\tabcolsep) * \real{0.3333}}
  >{\raggedright\arraybackslash}p{(\linewidth - 4\tabcolsep) * \real{0.3333}}@{}}
\toprule\noalign{}
\begin{minipage}[b]{\linewidth}\raggedright
Comparator
\end{minipage} & \begin{minipage}[b]{\linewidth}\raggedright
Purpose
\end{minipage} & \begin{minipage}[b]{\linewidth}\raggedright
Permission that must accompany each result
\end{minipage} \\
\midrule\noalign{}
\endhead
\bottomrule\noalign{}
\endlastfoot
Native LEFT / causal endpoint LEFT & Strong joint time--frequency
reference & Target training, scaler, selected checkpoint, endpoint
versus retrospective readout \\
TimeRadar & Pretrained time--frequency reference & Official weight
identity, pretraining exposure, target statistics, native offline versus
causal permission \\
History Mean 512 & Strong simple temporal control & Strictly earlier 512
observations; its own legal calibration and fixed residual scale \\
History and conditional Ridge & TEP history/conditional-response
controls & Normal FIT runs, lag specification, SELECT rule, legal
conditions, and own CAL \\
Fixed simple LEFT--Ridge combination & Tests whether joint modeling
exceeds simple evidence combination & Matching axes and calibration
roles; previously fixed rule, without QUERY weight selection \\
\end{longtable}

Existing raw results are reused only when data identity, channel order,
units, score axis, weight identity, and calibration role are compatible.
A missing compatible row is marked unavailable; an author-reported best
score under a different protocol is not inserted as a matched result.
This draft does not require retraining completed external models.

The candidate's fixed internal comparisons are the full
alternating-domain model, a parameter-matched time-only version without
the Haar-coordinate blocks, and a local-context version with a shorter
temporal receptive field. A fixed identity-transform Student-t score
tests whether learned transformations are necessary. These controls
distinguish learned density from the prior alone, multiscale coordinates
from temporal capacity, and long context from a local likelihood model.
They do not individually prove weak-anomaly or lag-change detection.
Known controlled tasks, reported separately from public-data scores, are
required for those claims. Stopped reference-center, masked-inference,
SOC, routing, and shared-prediction candidates are not resumed.

Before opening the new QUERY results, we also fix the zero-parameter
untrained full-model score as a control. All affine couplings and the
input affine start at identity, leaving one net orthogonal Haar analysis
followed by the same Student-t prior. This control is computed directly
from observed windows without fitting or additional training; it
separates learning from the change in an independent non-Gaussian prior
under coordinate rotation. It is not an alternative score chosen from
QUERY performance.

\subsection{Calibration, events, and sustained
detection}\label{calibration-events-and-sustained-detection}

Each method obtains a threshold from its own score units using the same
authorized CAL locations and the frozen quantile rule. The existing rule
is the 0.99 higher empirical quantile with a strict
score-greater-than-threshold decision. A score-definition change
therefore requires its own CAL evaluation; a historical numerical
threshold is not transferable across units. CAL is not a QUERY-derived
false-alarm guarantee, and its finite size and domain shift are
disclosed. No QUERY labels choose a threshold, checkpoint, seed,
horizon, aggregation weight, or entity.

TEP uses all 8,000 authorized normal CAL endpoints. PSM and SWaT use the
same 512 pre-fixed mature CAL endpoints as their cached native controls.
The new SWaT model can legally score earlier CAL endpoints with its
shorter history, but the comparison retains the native control's
original mature axis beginning at 446,173; the extra predictions are
archived rather than substituted into calibration.

Pointwise Average Precision (AP) and ROC-AUC are computed on the exact
common observed axis. AP is not renamed trapezoidal PR-AUC. Strict
pointwise precision, recall, and F1 use the frozen calibrated threshold,
without point adjustment. Three-consecutive-point confirmation is
causal: confirmation occurs at the third scored positive point, without
retrospectively marking the first two as detected. Both confirmed event
discovery and confirmed anomalous-point coverage are reported.

Event evaluation uses the original event boundaries and a fixed
one-to-one matching rule between predicted alarm fragments and true
events. Precision, recall, and F1 under that matching are distinct from
any-hit recall, which is reported separately. We report the first
detection delay, undetected-event censoring, normal false-positive
points, and false-alarm fragments per 1,000 observed normal points.
Within-event coverage, predefined long-event steady-state macro and
pooled coverage, and the longest consecutive miss quantify persistence
rather than onset alone. Physical-time delay is supplied only when the
sample interval is reliable. Unscored points are unknown, not
automatically normal or missed. Discontinuous query segments are
evaluated with their boundaries intact.

\subsection{Reproducibility, uncertainty, and decision
criteria}\label{reproducibility-uncertainty-and-decision-criteria}

The initial complete candidate and matched controls use the fixed first
seed 2025. Configurations record source and data versions, role
manifests, actual training updates, selection decisions, parameter
counts, runtime, peak memory, and failure costs. Complete continuous
scores, original timestamps or indices, labels for evaluation, and
calibration inputs are retained. Numerical checks verify inverse
reconstruction, the analytic log-determinant on a minimal input, finite
gradients, causal endpoint visibility, and equality between the
implemented training criterion and evaluation score. These are validity
checks, not evidence of detection improvement.

TEP and PSM reuse the frozen normal training draw plans and maximum
update counts of the cached target-trained comparison (3,760 and 12,400
updates). SWaT uses the same maximum of 20 epochs and normal-SELECT
patience of three; actual stopping and selection can differ between
objectives. The new SWaT implementation halves its learning rate after
each epoch. The historical native implementation's type-1 schedule may
have a one-epoch phase difference, so this row is not described as an
identical per-update optimizer trajectory. The three new learned
variants share their actual schedule and selection rule exactly.

A single seed on exposed development data supports screening, not a
statistically confirmed generalization claim. Event-level paired
summaries retain every event, while TEP uncertainty uses simulation run
as a clustering unit; overlapping points are not independent
replications. If the candidate demonstrates useful detection, subsequent
confirmation freezes the implementation before new independent entities
or runs are opened and includes repeated seeds. No significance claim is
made from an event count that overstates the number of independent
operating runs.

The candidate is supported only if an improvement relevant to weak or
sustained anomaly detection can be attributed to the learning criterion
or its temporal/multiscale structure and is not purchased through a
severe increase in normal false alarms. It need not dominate every
dataset or every metric. Conversely, a small AP change alone, an overall
likelihood rescaling, improved training loss, or a threshold-only gain
does not establish a methodological contribution. If the fixed complete
design lacks a useful effect relative to the mature methods and simple
controls, this candidate is stopped without a sequence of post hoc
score, threshold, or architectural adjustments.

\bibliography{references.bib}

\end{document}
